\documentclass[conference]{IEEEtran}
\IEEEoverridecommandlockouts

\usepackage{cite}
\usepackage{amsmath,amssymb,amsfonts}
\usepackage{algorithmic}
\usepackage{algorithm}
\usepackage{graphicx}
\usepackage{textcomp}
\usepackage{xcolor}
\usepackage{booktabs}
\usepackage{url}
\usepackage{microtype}

\def\BibTeX{{\rm B\kern-.05em{\sc i\kern-.025em b}\kern-.08em
    T\kern-.1667em\lower.7ex\hbox{E}\kern-.125emX}}

\begin{document}

\title{Bridging Diffusion Models and Flow Matching: \\ A Controlled Empirical Investigation and Few-Step Trajectory Straightening}

\author{\IEEEauthorblockN{Aditya Vardhan Singh}
\IEEEauthorblockA{\textit{Department of Computer Science} \\
\textit{Independent Research}}
}

\maketitle

\begin{abstract}
Continuous generative modeling has undergone a paradigm shift from traditional Denoising Diffusion Probabilistic Models (DDPM) toward Continuous Normalizing Flows and Conditional Flow Matching (CFM). While diffusion models conceptualize generation as reversing a Markovian noise-perturbation process, flow matching formulates synthesis as deterministic probability transport governed by continuous-time ordinary differential equations (ODEs). This paper provides a comprehensive, controlled empirical investigation comparing both methodologies under strictly aligned experimental controls: an identical 3.45M parameter U-Net architecture, matched training datasets (MNIST at $32 \times 32$ resolution), identical optimization hyperparameters, and hardware execution on Apple Silicon GPU via Metal Performance Shaders (\texttt{mps}). We analyze the models before and after bridging their theoretical formulations. Before combining, empirical evaluations demonstrate that while Flow Matching achieves structured sample contours at lower function-evaluation budgets ($NFE = 10$) due to linear optimal-transport paths, deterministic DDIM requires higher step counts or stochastic Langevin noise injection ($\eta = 1.0$) to eliminate high-frequency background drift. Following this baseline study, we bridge the two paradigms using 2-Rectified Flow (Reflow) to decouple and straighten state-space transport trajectories. By retraining on deterministically coupled endpoints $(Z_0, Z_1)$ generated from the baseline flow, 2-Rectified Flow eliminates trajectory crossings, enabling crisp, full-contrast image generation in a single Euler integration step ($NFE = 1, \sim 36$\,ms, $\text{max} = 1.000, \sigma = 0.245$). We discuss the computational trade-offs, theoretical implications, and scaling pathways for modern generative pipelines.
\end{abstract}

\begin{IEEEkeywords}
Generative models, diffusion probabilistic models, conditional flow matching, rectified flow, ordinary differential equations, trajectory straightening, one-step generation.
\end{IEEEkeywords}

\section{Introduction}
\IEEEPARstart{G}{enerative} modeling of continuous, high-dimensional distributions has become the cornerstone of modern visual synthesis. Over the past several years, Denoising Diffusion Probabilistic Models (DDPM) \cite{ho2020denoising} and Score-Based Generative Models \cite{song2019generative, song2020score} have emerged as the preeminent framework, displacing Generative Adversarial Networks (GANs) \cite{goodfellow2014generative} due to their stable maximum-likelihood-based training dynamics and superior mode coverage. Seminal text-to-image systems such as DALL-E 2 \cite{ramesh2022hierarchical} and Stable Diffusion \cite{rombach2022high} rely heavily on this formulation.

Despite their expressive capacity, traditional diffusion models suffer from a fundamental computational limitation: inference requires iteratively stepping through a discretized reverse stochastic process or solving a stiff probability flow ordinary differential equation (ODE) \cite{song2020denoising}. This sampling procedure typically requires dozens to hundreds of neural network function evaluations (NFE), imposing substantial latency and computational overhead that impedes real-time deployment.

To circumvent this bottleneck, recent research has gravitated toward Continuous Normalizing Flows (CNFs) \cite{chen2018neural} through the lens of Conditional Flow Matching (CFM) \cite{lipman2022flow} and Rectified Flow \cite{liu2022flow, albergo2022building}. Rather than perturbing data along a fixed stochastic schedule, flow matching directly regresses a time-dependent vector field that transports samples from a simple base distribution (e.g., standard Gaussian noise) to the empirical data distribution along straight interpolation paths. This formulation underpins state-of-the-art foundation architectures, including Stable Diffusion 3 \cite{esser2024scaling} and FLUX.1.

Notwithstanding these advances, empirical comparisons in literature are frequently confounded by discrepancies in backbone parameter capacity, disparate training budgets, differing dataset normalizations, or conflation of the training loss with the numerical ODE solver. In this paper, we establish a controlled, head-to-head empirical investigation to answer three fundamental questions:
\begin{enumerate}
    \item \textbf{Controlled Optimization}: Under strictly identical backbone architectures, training budgets, and datasets, how do diffusion objectives ($\epsilon$-prediction) and flow matching objectives (velocity regression) compare in training stability and numerical convergence?
    \item \textbf{Compute vs. Quality Frontier}: How does synthesis fidelity scale as a function of sampling compute (NFE) across deterministic (DDIM, Euler) and stochastic (DDPM) solvers?
    \item \textbf{Trajectory Straightening}: Can the gap between multi-step diffusion and single-step synthesis be bridged through trajectory straightening (Reflow) without relying on unstable adversarial losses or complex distillation objectives?
\end{enumerate}

\section{Mathematical Framework and Unification}

\subsection{Variance-Preserving Diffusion Models}
Let $x_1 \sim p_{\text{data}}(x)$ denote a clean sample drawn from the data distribution on $\mathbb{R}^d$. The forward noising process is governed by a discrete Markov chain over $T$ timesteps with variance schedule $\beta_1, \dots, \beta_T \in (0, 1)$:
\begin{equation}
q(x_t | x_{t-1}) = \mathcal{N}\left(x_t; \sqrt{1 - \beta_t} x_{t-1}, \beta_t \mathbf{I}\right).
\end{equation}
Defining $\alpha_t = 1 - \beta_t$ and $\bar{\alpha}_t = \prod_{s=1}^t \alpha_s$, the marginal distribution at any arbitrary timestep $t$ admits a closed-form Gaussian representation:
\begin{equation}
q(x_t | x_1) = \mathcal{N}\left(x_t; \sqrt{\bar{\alpha}_t} x_1, (1 - \bar{\alpha}_t) \mathbf{I}\right),
\end{equation}
which enables generating noisy states via reparameterization:
\begin{equation}
x_t = \sqrt{\bar{\alpha}_t} x_1 + \sqrt{1 - \bar{\alpha}_t} \epsilon, \quad \epsilon \sim \mathcal{N}(0, \mathbf{I}).
\end{equation}

In the continuous-time limit ($T \to \infty$), this forward process corresponds to the It\^{o} stochastic differential equation (SDE):
\begin{equation}
dx_t = -\frac{1}{2} \beta(t) x_t dt + \sqrt{\beta(t)} dw_t,
\end{equation}
where $w_t$ is standard Brownian motion. Song et al. \cite{song2020score} demonstrated that this SDE possesses an identical marginal density trajectory to the deterministic \textit{probability flow ODE}:
\begin{equation}
\frac{dx_t}{dt} = -\frac{1}{2} \beta(t) x_t - \frac{1}{2} \beta(t) \nabla_{x_t} \log p_t(x_t).
\label{eq:prob_flow_ode}
\end{equation}
The score function $\nabla_{x_t} \log p_t(x_t)$ is linked to the noise prediction network $\epsilon_\theta(x_t, t)$ via:
\begin{equation}
\nabla_{x_t} \log p_t(x_t) = -\frac{\epsilon_\theta(x_t, t)}{\sqrt{1 - \bar{\alpha}_t}}.
\end{equation}
The model is trained via the reweighted variational lower bound (MSE):
\begin{equation}
\mathcal{L}_{\text{diff}}(\theta) = \mathbb{E}_{t \sim \mathcal{U}\{1, T\}, x_1, \epsilon} \left[ \|\epsilon_\theta(x_t, t) - \epsilon\|^2 \right].
\end{equation}

Inference can be executed via Denoising Diffusion Implicit Models (DDIM) \cite{song2020denoising} over a subsampled sequence of steps $\tau_1 < \dots < \tau_S$:
\begin{align}
x_{\tau_{i-1}} = &\sqrt{\bar{\alpha}_{\tau_{i-1}}} \left( \frac{x_{\tau_i} - \sqrt{1 - \bar{\alpha}_{\tau_i}} \epsilon_\theta(x_{\tau_i})}{\sqrt{\bar{\alpha}_{\tau_i}}} \right) \nonumber \\
&+ \sqrt{1 - \bar{\alpha}_{\tau_{i-1}} - \sigma_i^2} \epsilon_\theta(x_{\tau_i}) + \sigma_i \epsilon,
\end{align}
where $\sigma_i = \eta \sqrt{(1 - \bar{\alpha}_{\tau_{i-1}})/(1 - \bar{\alpha}_{\tau_i})} \sqrt{1 - \bar{\alpha}_{\tau_i}/\bar{\alpha}_{\tau_{i-1}}}$. Setting $\eta = 0$ yields a deterministic solver for the probability flow ODE (\ref{eq:prob_flow_ode}), whereas $\eta = 1$ recovers the ancestral stochastic DDPM sampler.

\subsection{Conditional Flow Matching and Rectified Flow}
Continuous Normalizing Flows model data generation via a time-dependent vector field $v_t: \mathbb{R}^d \to \mathbb{R}^d$ generating a flow $\phi_t$:
\begin{equation}
\frac{d}{dt} \phi_t(x) = v_t(\phi_t(x)), \quad \phi_0(x) = x.
\end{equation}
The density $p_t$ induced by pushing forward $p_0 = \mathcal{N}(0, \mathbf{I})$ satisfies the continuity equation:
\begin{equation}
\frac{\partial p_t}{\partial t} + \nabla \cdot (p_t v_t) = 0.
\end{equation}
While learning the marginal vector field $v_t$ directly is intractable, Lipman et al. \cite{lipman2022flow} and Liu et al. \cite{liu2022flow} proved that regressing a conditional vector field $u_t(x | x_1)$ yields an equivalent objective:
\begin{equation}
\begin{aligned}
&\mathbb{E}_{t, x_t} \left[ \| v_\theta(x_t, t) - v_t(x_t) \|^2 \right] \\
&\equiv \mathbb{E}_{t, x_0, x_1} \left[ \| v_\theta(x_t, t) - u_t(x_t | x_0, x_1) \|^2 \right].
\end{aligned}
\end{equation}

In Optimal Transport Conditional Flow Matching (OT-CFM) and 1-Rectified Flow, the probability path between noise $x_0 \sim \mathcal{N}(0, \mathbf{I})$ and data $x_1 \sim p_{\text{data}}$ is a straight-line interpolation:
\begin{equation}
x_t = (1 - t) x_0 + t x_1, \quad t \in [0, 1].
\end{equation}
The corresponding conditional velocity target is constant with respect to time along the trajectory:
\begin{equation}
u_t(x_t | x_0, x_1) = \frac{d x_t}{dt} = x_1 - x_0.
\end{equation}
The neural network $v_\theta(x, t)$ is trained via ordinary regression:
\begin{equation}
\mathcal{L}_{\text{CFM}}(\theta) = \mathbb{E}_{t, x_0, x_1} \left[ \| v_\theta(x_t, t) - (x_1 - x_0) \|^2 \right].
\end{equation}

\begin{algorithm}[t]
\caption{Training Baseline Models (Diffusion vs. Flow Matching)}
\label{alg:training_comparison}
\begin{algorithmic}[1]
\REQUIRE Dataset $p_{\text{data}}$, Backbone $f_\theta$, Method $\mathcal{M} \in \{\text{DIFF}, \text{CFM}\}$
\STATE Initialize model parameters $\theta$
\REPEAT
    \STATE Sample clean batch $x_1 \sim p_{\text{data}}$
    \IF{$\mathcal{M} = \text{DIFF}$}
        \STATE Sample timesteps $t \sim \mathcal{U}\{0, T-1\}$ and noise $\epsilon \sim \mathcal{N}(0, \mathbf{I})$
        \STATE Compute $x_t = \sqrt{\bar{\alpha}_t} x_1 + \sqrt{1 - \bar{\alpha}_t} \epsilon$
        \STATE Predict $\hat{\epsilon} = f_\theta(x_t, t/T)$
        \STATE $\mathcal{L} = \|\hat{\epsilon} - \epsilon\|^2$
    \ELSIF{$\mathcal{M} = \text{CFM}$}
        \STATE Sample continuous time $t \sim \mathcal{U}(0, 1)$ and base noise $x_0 \sim \mathcal{N}(0, \mathbf{I})$
        \STATE Compute $x_t = (1 - t) x_0 + t x_1$
        \STATE Predict velocity $\hat{v} = f_\theta(x_t, t)$
        \STATE $\mathcal{L} = \|\hat{v} - (x_1 - x_0)\|^2$
    \ENDIF
    \STATE Update $\theta \leftarrow \theta - \gamma \nabla_\theta \mathcal{L}$ via AdamW
\UNTIL{convergence}
\end{algorithmic}
\end{algorithm}

\section{System Architecture and Implementation}

\subsection{Backbone Neural Architecture}
To guarantee that empirical differences stem strictly from the generative objective and sampling mechanics, both models share an identical, custom-engineered \texttt{CompactUNet} architecture:
\begin{itemize}
    \item \textbf{Input/Output Layers}: Accepts tensors of shape $(B, C, H, W) = (B, 1, 32, 32)$. An initial $3 \times 3$ convolution projects the input to a base channel width of $C_0 = 32$.
    \item \textbf{Encoder (Downsampling)}: Three hierarchical stages with channel multipliers $(1, 2, 4)$, yielding channel capacities of 32, 64, and 128. Each stage contains two residual blocks (Conv-GroupNorm-SiLU) followed by a strided convolution downsampler ($4 \times 4$ kernel, stride 2, padding 1).
    \item \textbf{Bottleneck}: Two full residual blocks operating at spatial resolution $4 \times 4$ with 128 channels.
    \item \textbf{Decoder (Upsampling)}: Three symmetric stages employing transposed convolutions ($4 \times 4$ kernel, stride 2, padding 1), skip-connection concatenation from the corresponding encoder stage, and two residual refinement blocks.
    \item \textbf{Time Conditioning}: Timesteps $t \in [0, 1]$ are encoded via continuous sinusoidal positional embeddings ($d = 32$), passed through a two-layer MLP with SiLU activations, and projected to a 128-dimensional latent vector injected additively into every residual block.
\end{itemize}
The total parameter count is exactly \textbf{3,455,713 parameters} ($13.8$\,MB in single-precision floating point).

\subsection{Experimental Controls}
Experiments were executed on an isolated macOS environment running PyTorch 2.8.0 accelerated via Metal Performance Shaders (\texttt{mps}) on an Apple Silicon unified memory GPU. Both models were trained on a standardized subset of 6,400 MNIST training images normalized strictly to $[-1, 1]$. Training hyperparameters were matched identically:
\begin{itemize}
    \item Batch size: $B = 64$ (100 updates per epoch).
    \item Total training epochs: $E = 10$ ($1,000$ total gradient updates).
    \item Optimizer: AdamW ($\text{lr} = 2 \times 10^{-4}, \beta_1 = 0.9, \beta_2 = 0.999$, weight decay $10^{-4}$).
    \item Gradient clipping: Norm clipped at $1.0$.
    \item Seed control: Global random seed pinned to 42.
\end{itemize}

\begin{figure}[t]
\centering
\includegraphics[width=\columnwidth]{results/eval/loss_curves.png}
\caption{Training loss curves over 1,000 optimization updates on Apple Silicon GPU (\texttt{mps}). Left: Diffusion baseline ($\epsilon$-prediction MSE). Right: Flow Matching baseline (velocity regression MSE).}
\label{fig:loss_curves}
\end{figure}

\section{Part I: Findings Before Combining (Controlled Baseline Comparison)}

\subsection{Optimization Stability and Dynamics}
As shown in Fig.~\ref{fig:loss_curves}, both formulations demonstrated smooth, stable convergence:
\begin{itemize}
    \item \textbf{Diffusion Baseline}: Average epoch loss decreased from $0.1464$ (Epoch 1) to $0.0268$ (Epoch 10). The gradient norm remained tightly bounded throughout training, with zero non-finite loss events. Total wall-clock training time was \textbf{362.4\,s} ($\sim 36.2$\,s/epoch).
    \item \textbf{Flow Matching Baseline}: Average velocity MSE loss decreased from $0.4500$ (Epoch 1) to $0.1635$ (Epoch 10). Total wall-clock training time was \textbf{352.9\,s} ($\sim 35.3$\,s/epoch).
\end{itemize}
The training throughput was virtually indistinguishable ($27.6$ steps/s for CFM vs. $26.9$ steps/s for Diffusion), confirming that neither objective imposes meaningful computational overhead during backpropagation.

\begin{table}[h]
\caption{Inference Latency, Throughput, and Computational Cost Across NFE Budgets ($B=16$ on \texttt{mps})}
\label{tab:sampling_benchmarks}
\centering
\begin{tabular}{cccccc}
\toprule
\textbf{NFE} & \multicolumn{2}{c}{\textbf{Diffusion (DDIM)}} & \multicolumn{2}{c}{\textbf{Flow Matching (Euler)}} & \textbf{Speedup} \\
\cmidrule(lr){2-3} \cmidrule(lr){4-5}
& Latency & Throughput & Latency & Throughput & (CFM vs. DDIM) \\
\midrule
5  & 203.0\,ms & 78.8\,img/s & 184.0\,ms & 87.0\,img/s & $1.10\times$ \\
10 & 401.8\,ms & 39.8\,img/s & 333.1\,ms & 48.0\,img/s & $1.21\times$ \\
20 & 717.0\,ms & 22.3\,img/s & 707.5\,ms & 22.6\,img/s & $1.01\times$ \\
50 & 2115.5\,ms & 7.56\,img/s & 1677.1\,ms & 9.54\,img/s & $1.26\times$ \\
\bottomrule
\end{tabular}
\end{table}

\subsection{Sampling Performance and Latency Trade-Offs}
In Table~\ref{tab:sampling_benchmarks}, we evaluate batch generation ($B=16$) across function-evaluation budgets $NFE \in \{5, 10, 20, 50\}$.
\begin{itemize}
    \item Per-step latency scales linearly with NFE for both methods ($\sim 34\text{--}42$\,ms per evaluation step).
    \item Flow matching exhibits slightly higher throughput across all budgets because the Euler solver requires fewer intermediate tensor operations (coordinate addition $x_{t+\Delta t} = x_t + \Delta t \cdot v_t$) compared to DDIM, which evaluates alpha-schedule ratios, noise predictions, and variance scaling terms at every step.
\end{itemize}

\subsection{Qualitative Sample Evolution and Failure Modes}
Fig.~\ref{fig:comparative_grid} illustrates generation quality across matching NFE budgets:
\begin{enumerate}
    \item \textbf{Low-NFE Regime ($NFE \le 10$)}: Flow matching with Euler integration resolves coherent digit silhouettes and clear foreground/background separation as early as 10 steps. In contrast, deterministic DDIM ($\eta = 0$) exhibits persistent salt-and-pepper background haze. Under constrained training compute (1,000 updates), the noise prediction network $\epsilon_\theta$ retains non-zero estimation error across timesteps; during large deterministic DDIM jumps, this residual error accumulates uncorrected.
    \item \textbf{Stochastic Compensation}: When diffusion is sampled with stochastic Langevin noise ($\eta = 1.0$) at $NFE \ge 50$, the injected Brownian motion counterbalances the deterministic drift, clearing background artifacts completely and producing crisp, high-contrast strokes.
\end{enumerate}

\begin{figure*}[t]
\centering
\includegraphics[width=0.68\textwidth,keepaspectratio]{results/eval/comparative_grid.png}
\caption{Controlled qualitative comparison between Diffusion (DDIM, left) and Flow Matching (Euler, right) across sampling budgets $NFE \in \{5, 10, 20, 50\}$ under identical model capacity and training compute.}
\label{fig:comparative_grid}
\end{figure*}

\section{Part II: Findings After Combining (2-Rectified Flow / Reflow)}

Although 1-Rectified Flow outperforms diffusion in the low-step regime ($NFE = 10$), pushing inference to the theoretical limit of a single Euler step ($NFE = 1$) fails catastrophically, yielding a blurry, low-contrast smear.

\subsection{The Trajectory Crossing Problem}
In 1-Rectified Flow, noise $x_0 \sim \mathcal{N}(0, \mathbf{I})$ and data $x_1 \sim p_{\text{data}}$ are sampled independently. Although each individual trajectory $x_t = (1-t)x_0 + t x_1$ is a straight line, trajectories originating from different noise-data pairs intersect in state space.

Because two distinct trajectories can pass through the same point $x$ at time $t$ with different target velocities ($u_t^{(A)} \ne u_t^{(B)}$), the learned marginal vector field $v_\theta(x, t)$ must compromise by outputting the conditional expectation:
\begin{equation}
v_t(x) = \mathbb{E}[x_1 - x_0 | x_t = x].
\end{equation}
This spatial averaging introduces significant curvature into the marginal flow lines. Integrating this curved field with large step sizes ($\Delta t = 1.0$) causes severe discretization truncation error.

\subsection{The 2-Rectified Flow (Reflow) Algorithm}
To straighten the flow lines, we implement \textbf{2-Rectified Flow} \cite{liu2022flow}, summarized in Algorithm~\ref{alg:reflow}:
\begin{enumerate}
    \item \textbf{Deterministic Endpoint Coupling}: Using the trained 1-Rectified Flow model $v^1_\theta$, we sample 3,200 initial noise vectors $Z_0 \sim \mathcal{N}(0, \mathbf{I})$ and integrate them to terminal states $Z_1$ using a high-precision 50-step ODE solver:
    \begin{equation}
    Z_1 = \text{ODE\_Solve}\left(v^1_\theta, Z_0, t \in [0, 1], \text{steps}=50\right).
    \end{equation}
    This establishes a causal, deterministic coupling between each noise vector $Z_0$ and its destination sample $Z_1$.
    \item \textbf{Trajectory Straightening}: We train a student vector field $v^2_\phi$ on the straight lines connecting these coupled pairs:
    \begin{equation}
    \begin{aligned}
    \mathcal{L}_{\text{reflow}}(\phi) = \mathbb{E}_{(Z_0, Z_1), t} [&\| v^2_\phi((1-t)Z_0 + t Z_1, t) \\
    &- (Z_1 - Z_0) \|^2].
    \end{aligned}
    \end{equation}
\end{enumerate}

\begin{algorithm}[t]
\caption{2-Rectified Flow (Reflow) Trajectory Straightening}
\label{alg:reflow}
\begin{algorithmic}[1]
\REQUIRE Trained 1-RF model $v^1_\theta$, Number of pairs $N=3200$, Steps $S=50$
\STATE \textbf{Phase 1: Generate Coupled Pairs}
\FOR{$k = 1$ to $N/B$}
    \STATE Sample $Z_0 \sim \mathcal{N}(0, \mathbf{I})$
    \STATE Integrate $Z_1 = \text{Euler\_Integrate}(v^1_\theta, Z_0, S)$
    \STATE Store paired endpoints $(Z_0, Z_1)$ in dataset $\mathcal{D}_{\text{paired}}$
\ENDFOR
\STATE \textbf{Phase 2: Train Reflowed Student Model}
\STATE Initialize student model $v^2_\phi \leftarrow v^1_\theta$
\REPEAT
    \STATE Sample batch $(Z_0, Z_1) \sim \mathcal{D}_{\text{paired}}$ and $t \sim \mathcal{U}(0, 1)$
    \STATE Interpolate $Z_t = (1 - t) Z_0 + t Z_1$
    \STATE Target velocity $U = Z_1 - Z_0$
    \STATE Update $\phi \leftarrow \phi - \gamma \nabla_\phi \| v^2_\phi(Z_t, t) - U \|^2$
\UNTIL{convergence}
\end{algorithmic}
\end{algorithm}

\begin{figure*}[t]
\centering
\includegraphics[width=0.68\textwidth,keepaspectratio]{results/eval/reflow_minimal_steps_comparison.png}
\caption{Few-step synthesis comparison between 1-Rectified Flow (left column) and 2-Rectified Flow / Reflow (right column) across budgets $NFE \in \{1, 2, 4, 8\}$ under identical starting noise vectors.}
\label{fig:reflow_comparison}
\end{figure*}

\subsection{Theoretical Guarantee of Straightening}
Because $(Z_0, Z_1)$ are generated along the non-intersecting trajectories of an ODE, the linear paths $(1-t)Z_0 + tZ_1$ do not cross in state space (or cross with vastly lower frequency). Liu et al. \cite{liu2022flow} proved that the total transport cost strictly decreases:
\begin{equation}
\int_0^1 \mathbb{E}\left[ \| v^2_\phi(Z_t, t) \|^2 \right] dt \le \int_0^1 \mathbb{E}\left[ \| v^1_\theta(X_t, t) \|^2 \right] dt.
\end{equation}
As a direct consequence, the velocity field becomes near-constant along each path ($v^2_\phi(Z_t, t) \approx Z_1 - Z_0$). Under this condition, a single forward evaluation Euler step:
\begin{equation}
x_{\text{gen}} = Z_0 + 1.0 \cdot v^2_\phi(Z_0, 0)
\end{equation}
yields an exact solution to the transport equation.

\subsection{Empirical Validation of Single-Step Synthesis}
Generating the 3,200 coupled pairs required 330.8\,s on \texttt{mps}. Fine-tuning the 2-Rectified Flow student for 10 epochs converged in 177.8\,s, with paired loss declining from $0.0241 \to 0.0101$ (over an order of magnitude lower than 1-RF training loss, reflecting the near-constant target velocity).

We evaluated both 1-RF and 2-RF on identical starting noise across few-step budgets $NFE \in \{1, 2, 4, 8\}$. As shown in Fig.~\ref{fig:reflow_comparison} and Table~\ref{tab:pixel_stats}:
\begin{itemize}
    \item \textbf{1-Rectified Flow at 1 Step}: Fails to synthesize digits, producing a faint, washed-out smear with a maximum pixel intensity of only $0.729$ and standard deviation $\sigma = 0.152$.
    \item \textbf{2-Rectified Flow at 1 Step}: Achieves full dynamic range ($\text{max} = 1.000, \sigma = 0.245$), generating sharp, high-contrast, fully rendered digits with deep black backgrounds in **one single forward pass ($\sim 36$\,ms)**.
\end{itemize}

\begin{table}[h]
\caption{Statistical Distribution and Contrast Across Few-Step Budgets}
\label{tab:pixel_stats}
\centering
\begin{tabular}{lcccc}
\toprule
\textbf{Model / Evaluation Budget} & \textbf{Min} & \textbf{Max} & \textbf{Mean} & \textbf{Std ($\sigma$)} \\
\midrule
1-Rectified Flow ($NFE = 1$) & 0.000 & 0.729 & 0.113 & 0.152 \\
\textbf{2-Rectified Flow ($NFE = 1$)} & \textbf{0.000} & \textbf{1.000} & \textbf{0.123} & \textbf{0.245} \\
1-Rectified Flow ($NFE = 2$) & 0.000 & 0.976 & 0.103 & 0.185 \\
\textbf{2-Rectified Flow ($NFE = 2$)} & \textbf{0.000} & \textbf{1.000} & \textbf{0.121} & \textbf{0.247} \\
1-Rectified Flow ($NFE = 4$) & 0.000 & 1.000 & 0.110 & 0.223 \\
\textbf{2-Rectified Flow ($NFE = 4$)} & \textbf{0.000} & \textbf{1.000} & \textbf{0.120} & \textbf{0.250} \\
\bottomrule
\end{tabular}
\end{table}

\section{Discussion and Limitations}
While Reflow successfully achieves 1-step synthesis without adversarial discriminators, several engineering trade-offs warrant consideration:
\begin{itemize}
    \item \textbf{Data Generation Overhead}: Coupling endpoints requires an offline ODE generation phase. In large-scale foundation models, integrating billions of latent pairs demands non-trivial compute clusters.
    \item \textbf{Dataset Complexity}: This study evaluated $32 \times 32$ grayscale imagery. In complex multimodal domains (e.g., ImageNet or high-resolution text-to-image), manifold geometry is vastly more intricate, requiring progressive distillation \cite{salimans2022progressive} or consistency training \cite{song2023consistency} alongside Reflow.
\end{itemize}

\section{Conclusion}
This study presented a controlled empirical comparison between traditional diffusion models and conditional flow matching under an identical 3.45M parameter backbone and fixed compute budget. We demonstrated that flow matching offers superior sample efficiency at moderate step counts ($NFE = 10\text{--}20$) compared to deterministic DDIM. Furthermore, by bridging the paradigms through 2-Rectified Flow trajectory straightening, we eliminated the multi-step sampling constraint entirely, enabling high-fidelity 1-step image generation in $\sim 36$\,ms.

\begin{thebibliography}{00}
\bibitem{ho2020denoising} J. Ho, A. Jain, and P. Abbeel, ``Denoising diffusion probabilistic models,'' in \textit{Proc. Adv. Neural Inf. Process. Syst. (NeurIPS)}, vol. 33, pp. 6840--6851, 2020.
\bibitem{song2019generative} Y. Song and S. Ermon, ``Generative modeling by estimating gradients of the data distribution,'' in \textit{Proc. Adv. Neural Inf. Process. Syst. (NeurIPS)}, vol. 32, 2019.
\bibitem{song2020score} Y. Song, J. Sohl-Dickstein, D. P. Kingma, A. Kumar, S. Ermon, and B. Poole, ``Score-based generative modeling through stochastic differential equations,'' in \textit{Proc. Int. Conf. Learn. Represent. (ICLR)}, 2021.
\bibitem{song2020denoising} J. Song, C. Meng, and S. Ermon, ``Denoising diffusion implicit models,'' in \textit{Proc. Int. Conf. Learn. Represent. (ICLR)}, 2021.
\bibitem{lipman2022flow} Y. Lipman, R. T. Q. Chen, H. Ben-Hamu, M. Nicklas, and M. Le, ``Flow matching for generative modeling,'' in \textit{Proc. Int. Conf. Learn. Represent. (ICLR)}, 2023.
\bibitem{liu2022flow} X. Liu, C. Gong, and Q. Liu, ``Flow straight and fast: Learning to generate and transfer data with rectified flow,'' in \textit{Proc. Int. Conf. Learn. Represent. (ICLR)}, 2023.
\bibitem{albergo2022building} M. S. Albergo and E. Vanden-Eijnden, ``Building normalizing flows with stochastic interpolants,'' in \textit{Proc. Int. Conf. Learn. Represent. (ICLR)}, 2023.
\bibitem{chen2018neural} R. T. Q. Chen, Y. Rubanova, J. Bettencourt, and D. K. Duvenaud, ``Neural ordinary differential equations,'' in \textit{Proc. Adv. Neural Inf. Process. Syst. (NeurIPS)}, vol. 31, 2018.
\bibitem{goodfellow2014generative} I. Goodfellow et al., ``Generative adversarial nets,'' in \textit{Proc. Adv. Neural Inf. Process. Syst. (NeurIPS)}, vol. 27, 2014.
\bibitem{rombach2022high} R. Rombach, A. Blattmann, D. Lorenz, P. Esser, and B. Ommer, ``High-resolution image synthesis with latent diffusion models,'' in \textit{Proc. IEEE/CVF Conf. Comput. Vis. Pattern Recognit. (CVPR)}, pp. 10684--10695, 2022.
\bibitem{ramesh2022hierarchical} A. Ramesh, P. Dhariwal, A. Nichol, C. Chu, and M. Chen, ``Hierarchical text-conditional image generation with CLIP latents,'' \textit{arXiv preprint arXiv:2204.06125}, 2022.
\bibitem{esser2024scaling} P. Esser et al., ``Scaling rectified flow transformers for high-resolution image synthesis,'' in \textit{Proc. Int. Conf. Mach. Learn. (ICML)}, 2024.
\bibitem{salimans2022progressive} T. Salimans and J. Ho, ``Progressive distillation for fast sampling of diffusion models,'' in \textit{Proc. Int. Conf. Learn. Represent. (ICLR)}, 2022.
\bibitem{song2023consistency} Y. Song, P. Dhariwal, M. Chen, and I. Sutskever, ``Consistency models,'' in \textit{Proc. Int. Conf. Mach. Learn. (ICML)}, pp. 32211--32252, 2023.
\end{thebibliography}

\end{document}
